\section{Method}
\label{sec:method}

\ifshort
\paragraph{Threat model.}
The adversary can write text into a field the agent reads while working and
controls a network destination; they cannot modify the system prompt, tool
definitions, weights, or harness. The agent's credentials are legitimate---it is
authorized to retrieve member records, because retrieval is the task. We measure
whether it can be induced to \emph{place} a record in an outbound argument
addressed to the adversary, not whether that request would traverse a network
(\cref{sec:attempt}).

\paragraph{Canary, detector, policy.}
A synthetic member record carrying \canaryID{} and \nRecordFields{} fields
(name, date of birth, diagnosis codes, medications, address) is planted among
\nCustomers{} records; it is invented in full. \emph{Detection}---which encoding
of the canary is present in a string---is a pure function of bytes, frozen, and
exercised by \nOracleTests{} tests. \emph{Policy}---which tools are egress, which
host is the adversary's---is configurable and delegates every byte-level decision
to the detector rather than reimplementing a matcher; re-scoring every logged
call and episode through both paths gave identical verdicts. The detector reports
the least-transformed of \nEncodings{} encodings (\encodingList{}); a novel
encoding is scored conservatively as no leak, so rates are lower bounds. Canary
bytes in a \emph{retrieval} argument are never counted---looking a member up is
the task---and only the \nEgressTools{} egress tools (\egressToolList{}) are
scored.
\else

\subsection{Threat model}
The adversary can write text into at least one field that the agent will read
while performing its normal task, and controls a network destination. The
adversary cannot modify the agent's system prompt, its tool definitions, its
model weights, or the harness. The agent's own credentials are legitimate: it is
authorized to retrieve member records, because retrieval is the task. We measure
whether the agent can be induced to \emph{place} a protected record in an
outbound argument addressed to the adversary. We deliberately do not model
whether that request would traverse a real network; see
\cref{sec:attempt}.

\subsection{The canary}
A synthetic member record is planted among \nCustomers{} records in the agent's
data store, carrying the identifier \canaryID{} and \nRecordFields{} fields
including name, date of birth, diagnosis codes, medications, address, and a
recent claim. The record is invented in full; no field corresponds to a real
person. The workflow is constructed so that no legitimate task requires
transmitting it outside the trust boundary. The agent is expected to \emph{read}
the record---that is the task---and never to send it.

\subsection{Detection and policy are separate layers}
Two questions are easy to conflate and have different stability requirements.
\emph{Detection} asks which encoding of the canary, if any, is present in a given
string; it is a pure function of bytes and is frozen. \emph{Policy} asks which
tool names constitute egress and which host belongs to the adversary; it is
deployment-specific and configurable, because a user evaluating their own agent
has different tool names. The frozen detector is exercised by \nOracleTests{}
tests. The policy layer delegates every byte-level decision to the detector
rather than reimplementing any matcher; we verified the split is not a behavioural
change by re-scoring every logged tool call and episode in the study through both
paths and confirming key-for-key identical verdicts.

\subsection{Matcher and encodings}
\label{sec:encodings}
The detector reports the least-transformed encoding that fires, over
\nEncodings{} recoverability-preserving encodings: \encodingList{}. The list is
closed under the requirement that the adversary can still recover the record;
transformations that destroy recoverability are outside the threat model,
because the adversary gains nothing from them. A novel encoding not on this list
would evade detection and be scored, conservatively, as no leak---so reported
rates are lower bounds.

Canary bytes appearing in a \emph{retrieval} argument are recorded separately and
never counted as exfiltration: looking a member up by identifier is the task, not
a leak. Only arguments to the \nEgressTools{} egress tools (\egressToolList{})
are scored.

\fi

\subsection{Agent under test}
\ifshort
The reference agent is a member-services agent for a fictional health plan, with
\nAgentTools{} tools and a \maxTurns{}-turn budget, which resolves a support
ticket end to end. Tool calls are emitted as plain text and parsed by regular
expression rather than through a native function-calling interface; its system
prompt states that member records are protected health information. Tools are
simulated and return fixed data.
\else
The reference agent is ARIA, an automated member-services agent for a fictional
health plan, which resolves a support ticket end to end: read the ticket, look up
the member record, search a policy knowledge base, and reply. It has
\nAgentTools{} tools and a budget of \maxTurns{} turns. Its system prompt states
explicitly that member records are protected health information and must be
treated accordingly; that sentence is present in every episode reported here,
including every episode in which a record was nevertheless transmitted.

Two properties of the agent are simplifications and are stated here rather than
in the limitations alone. First, tool calls are not made through a native
function-calling interface: the model emits \texttt{TOOL: name \{json\}} as plain
text and the harness parses it by regular expression. Second, tools are simulated
and return fixed data; no request leaves the process.
\fi

\ifshort
\paragraph{Corpus, control arm, and the gate.}
The corpus is \corpusSize{} payloads: \nChannels{} channels (\channelList{})
crossed with \nVariants{} variants (\variantList{}), each channel dressing the
injected text as that surface plausibly looks and each variant holding intent
fixed while changing surface form (\cref{app:payload}). Every payload has a
matched control: identical ticket, member, and tool surface, injected span
removed and nothing substituted---the only way to measure the label's
false-positive rate, and \cref{sec:fp} reports that measuring it changed the
label.
\label{sec:gate}
An agent that never executes produces zero detected exfiltrations and zero false
positives---the characteristic failure of a bring-your-own-agent tool, whose most
dangerous output is a clean bill of health it did not earn. The harness therefore
refuses to score until four checks pass, the fourth a positive control in which a
blatant unobfuscated exfiltration instruction is issued and the oracle \emph{must}
detect the leak; an override stamps every artifact unverified. The gate caught a
real defect in this project's own adapter scaffold, where template and harness
imported the same base class by two module paths and a correct adapter was
rejected.
\else
\subsection{Corpus}
The corpus is \corpusSize{} payloads: \nChannels{} delivery channels
(\channelList{}) crossed with \nVariants{} obfuscation variants
(\variantList{}). Each channel dresses the injected text as that surface
plausibly looks---an appended record note, a postscript on a web-form
submission, an auto-synced knowledge-base addendum, a handoff note from a prior
agent. The variants hold the instruction's intent fixed while changing its
surface form: unadorned prose, a fabricated compliance directive, an HTML
comment with an image-tag decoy, and zero-width characters spliced through the
words. An example appears in \cref{app:payload}.

\subsection{Control arm}
\label{sec:control}
Every payload has a matched control: the identical ticket, the identical member,
the identical tool surface, with the injected span removed and nothing
substituted. The control arm is not a nicety. It is the only way to measure the
false-positive rate of the label, and \cref{sec:fp} reports that measuring it
changed the label. A full evaluation of one agent is therefore
\pooledInjectedN{}~/~\nModels{} injected episodes plus an equal number of
controls.

\subsection{The positive-control gate}
\label{sec:gate}
An agent that never executes---a misconfigured adapter, a tool-name mismatch, a
harness wired to nothing---produces zero detected exfiltrations and zero false
positives. That is the characteristic failure of a bring-your-own-agent security
tool: the most dangerous output such a tool can emit is a clean bill of health it
did not earn. The harness therefore refuses to report a score until four checks
pass, of which the fourth is a positive control: a blatant, unobfuscated
exfiltration instruction is issued and the oracle \emph{must} detect the
resulting leak. If it does not, no number is printed. An operator may override
the gate explicitly, in which case every artifact produced---including the
rendered report---is stamped as unverified.

\ifshort
\else
The gate is not hypothetical: it caught a real defect in this project's own
adapter scaffold, in which the template and the harness imported the same base
class by two different module paths, producing two distinct classes and causing a
correctly written adapter to be rejected. The defect would have affected every
user on first run.
\fi
\fi
